\chapter{Line Fitting-Based Corner-Like Detector for 2D Laser Scanners Data}\label{chap:features}

TBC

\section{Background}\label{sec:features:background}

TBC (background / literature analysis specific to the topic of this chapter)

\subsection{Line fitting intersection}\label{sec:features:background:line-fitting}

A trend in the literature is to extract corners at the intersections of
line segments detected in the 2D~laser
scan~\cite{castellanos:mfi:1996,roumeliotis-bekey:iros:2000,jung:tits:2008,falomir:prl:2013,fernandez-moral:ijrr:2015,mohamed:isprs-arxiv:2016,zhao:icsess:2016,lin:yac:2017,qu:robio:2018,liu:electronics:2022,liu:sensors:2023}.
First, the laser data may be filtered to remove invalid values
(zero range values~\cite{jung:tits:2008}) and possible noise
(\eg, considering the minimum distance to the two neighbors of
a certain point to remove isolated data~\cite{jung:tits:2008},
mean average filtering~\cite{liu:electronics:2022}, or bilateral filters
to remove non-linear noise~\cite{liu:sensors:2023}).
Next, a segmentation step extracts sets of continuous points
in the environment for line fitting estimation.
Castellanos~\etal~\cite{castellanos:mfi:1996} segment the scan into
polygonal lines and subsequently divide them into
segments using an iterative method.
Falomir~\etal~\cite{falomir:prl:2013} search for discontinuities
in the data to form point subsets and use a recursive line-fitting
algorithm~\cite{xavier:icra:2005} to obtain line segments.
Fernández-Moral~\etal~\cite{fernandez-moral:ijrr:2015} search for
parameters of a 2D~line that maximize the number of points
supporting the line model using the \gls{ransac}
algorithm~\cite{fischler-bolles:cacm:1981}.
Mohamed~\etal~\cite{mohamed:isprs-arxiv:2016} use the line tracking
algorithm~\cite{munoz:ropec:2013} to iteratively segment lines
in the scan and perform \gls{pca} to compute the lines' parameters.
Then, the intersection points between consecutive line segments
define a corner feature in the environment.

Furthermore, Lin~\etal~\cite{lin:yac:2017} and
Liu~\etal~\cite{liu:electronics:2022,liu:sensors:2023} use adaptive threshold
segmentation to cluster range data into continuous spatial subsets.
Although Lin~\etal~\cite{lin:yac:2017} state that an adaptive segmentation
was used due to the non-uniform distribution of the 2D~laser data,
the authors do not specify the threshold used in their work.
In Liu~\etal~\cite{liu:electronics:2022} work, the threshold is based on
the laser's angular resolution and the distance to
the current point under analysis.
Liu~\etal~\cite{liu:sensors:2023} also establish a threshold depending on
the angular resolution. Instead of considering only the distance of
the point under analysis, the proposed approach~\cite{liu:sensors:2023} selects
the larger distance between adjacent points.
Even though Roumeliotis and Bekey~\cite{roumeliotis-bekey:iros:2000} do not use
a specific adaptive threshold, the proposed SEGMENTS algorithm with two-level
\gls{ekf} exhibits adaptive behavior. A stricter filter,
in terms of the system's noise covariance, attempts to follow a straight line,
while a more flexible \gls{ekf} tries to follow a new segment during each scan.
The four methods~\cite{lin:yac:2017,liu:electronics:2022,liu:sensors:2023,roumeliotis-bekey:iros:2000}
also define a corner feature as the intersection
between consecutive line segments.

A specific application of corner detection with 2D~laser data in the literature
is vehicle detection~\cite{jung:tits:2008,qu:robio:2018}.
Jung~\etal~\cite{jung:tits:2008} cluster laser data using a fixed threshold
between consecutive range values. Then, the method performs template matching
for every point under analysis, assuming the point is
the vertex of the L-shape of a vehicle in the scan.
Jung~\etal~\cite{jung:tits:2008} use \gls{svd} to estimate the L-shape's
line parameters and decompose the problem into line-curve estimation
when the vehicle has a curved corner in the data.
In contrast, Qu~\etal~\cite{qu:robio:2018} segment the data using
\gls{dbscan}~\cite{ester:kdd:1996}.
The method performs L-shape fitting in two steps:
searching for the L-shape vertices and corner-point localization,
the latter based on \gls{ransac}~\cite{fischler-bolles:cacm:1981}.

However, corner detection algorithms that rely on global search during
line segment scanning may be limited to well-structured environments.
Thus, we propose a corner-like detection method based on local line fitting,
without requiring global line detection or recursive data search.
Also, the experiments presented in \secref{sec:features:results} demonstrate
the ability of the proposed detector to extract more points of interest
from the scan in addition to environment corners.

\subsection{Curvature-based detection}\label{sec:features:background:curvature}

Other approaches to detecting corners in 2D~laser data estimate the curvature.
N\'{u}\~{n}ez~\etal~\cite{nunez:icra:2006} propose an adaptive curvature
estimator that compares the $k$-nearest points on
both sides of the point being evaluated.
The $k$ size is adjusted based on the distance between possible corners.
A segmentation step precedes the curvature estimation to detect breakpoints
using the adaptive breakpoint detector proposed by
Borges~and~Aldon~\cite{borges-aldon:jint:2004}.
Then, N\'{u}\~{n}ez~\etal~\cite{nunez:icra:2006} extract corner, line, and
curve segment features from the scan data by analyzing whether
there are peaks in curvature, flat segments with zero average curvature,
and flat segments with non-zero average curvature, respectively.
V\'{a}zquez-Mart\'{i}n~\etal~\cite{vazquez-martin:sensors:2009} use
the adaptive curvature estimator of N\'{u}\~{n}ez~\etal~\cite{nunez:icra:2006}
to segment the scan into clusters of homogeneous curvature.
Even though the existence of a corner is determined from the curvature function,
its keypoint definition is based on the intersection of the two line features
that generate the corner, to improve robustness to noise.
Also, Certad~\etal~\cite{certad:andescon:2012} modify the work of
N\'{u}\~{n}ez~\etal~\cite{nunez:icra:2006} to consider
two curvature estimations based on different noise-related thresholds,
thereby improving the rate of true detection for
corner and edge features in the experiments.

Moreover, Feng~\etal~\cite{feng:aici:2009} and
Zhao~\etal~\cite{zhao:icsess:2016} report a correlation between
the angle difference between the two fitted lines on
both sides of a scan point and the curvature value.
Feng~\etal~\cite{feng:aici:2009} introduce the micro-tangent line function
as the variance between the tangents at two consecutive points in the scan.
Similarly to N\'{u}\~{n}ez~\etal~\cite{nunez:icra:2006},
line and curve segments are extracted from consecutive range readings
with variance close to zero and with constant value larger than a threshold,
respectively, and corner features are extracted from peaks in the variance
function. Instead of only considering two consecutive points,
Zhao~\etal~\cite{zhao:icsess:2016} define two distance-based regions on
both sides of a point to estimate direction vectors and,
consequently, the angle difference between those vectors.
Candidate corner points are selected based on a fixed threshold
for the angle difference. Then, the corner point is estimated
via non-minimal suppression.

Still, most curvature-based methods for corner feature extraction associate
the keypoint with a laser reading. Given the discrete, radial nature
of 2D~laser, the estimated corner may vary farther away.
Also, there is the question of whether the curvature estimators
are invariant to scale, angular resolution, and different viewpoints.
We propose a detector based on the stability of intersection points
from fitted lines at different distance-based scales.
The experiments in \secref{sec:features:results} focus
on evaluating the invariance of the proposed corner-like features.

\section{Corner-Like Feature Detector}\label{sec:features:method}

The proposed methodology for feature extraction on 2D~laser data performs
local line fitting to detect corner-like points in the scan.
In summary, first, the scan data is transformed from polar to Cartesian
coordinates. This transformation maintains the data's angular coordinate order.
For each point in the scan, multiple distance-based scales are used
to extract arrival and departure neighborhoods,
depending on whether the latter are in the negative or positive directions
relative to the point's angular coordinate, respectively.
Next, a linear regression algorithm fits lines to each neighborhood.
The intersection point between the arrival and departure neighborhoods
is a candidate feature keypoint at a given scale.
Then, the following indicators are estimated to determine
if the candidate keypoint is a corner-line feature or not:
minimum number of points to fit a line,
angle between arrival and departure lines,
maximum \gls{rmse} allowed for the line fitting procedure, maximum distance
of the candidate point to the intersection centroid on the multiple scales,
and maximum distance to the laser point on analysis.
The candidate feature keypoint is validated if it passes
all previous indicators over a minimum number of consecutive scales,
thereby improving robustness to noise in the laser data.

\subsection{Linear least-squares regression}\label{sec:features:method:linear-regression}

A generic definition for the 2D~line model is presented in
\eqref{eq:feature:line-models:generic}, where $\left(x,y\right)$ are
the 2D~point coordinates and $\gamma_0$, $\gamma_1$, and $\gamma_2$
are the model parameters. However, the solution for the generic definition
estimated with a linear least-squares formulation
from a set of 2D~points would be a null vector.
\begin{equation}\label{eq:feature:line-models:generic}
\gamma_0 + \gamma_1 x + \gamma_2 y = 0
\end{equation}

An alternative formulation is to characterize the 2D~line relative to one
of the point coordinates, reducing from three to two parameters in the model.
The 2D~line formulation illustrated in \eqref{eq:feature:line-models:ymxb}
is relative to the $y$ coordinate with respect to $x$, where $\beta_1$ is
the line slope and $\beta_0$ is the line intersection on the y-axis.
The least-squares solution for the parameters $\beta_1$ and $\beta_0$
from a set of points $\left(x_i,y_i\right),i=1,\dots,N$ is formulated
in \eqref{eq:feature:line-models:ymxb:m} and
\eqref{eq:feature:line-models:ymxb:b}, where $\left(\bar{x},\bar{y}\right)$ is
the centroid point of the data~\cite{kenney-keeping:book:1962}.
Still, this model would not be valid when
$\sum_{i=1}^{n} \left(x_i - \bar{x}\right)^2=0$,
given the indeterminate form caused in the parameter $\hat{\beta}_1$ (\eg,
all the data points with the same $x$ coordinate and different $y$ values).
\begin{equation}\label{eq:feature:line-models:ymxb}
y = \beta_0 + \beta_1 x
\end{equation}
\begin{equation}\label{eq:feature:line-models:ymxb:m}
\hat{\beta}_1 = \frac{\sum_{i=1}^{n} \left(x_i - \bar{x}\right)\left(y_i - \bar{y}\right)}{\sum_{i=1}^{n} \left(x_i - \bar{x}\right)^2}
\end{equation}
\begin{equation}\label{eq:feature:line-models:ymxb:b}
\hat{\beta}_0 = \bar{y} - \hat{\beta}_1 \bar{x}
\end{equation}

The 2D~line can be defined relative to $x$ with respect to the $y$ coordinate,
as formulated in \eqref{eq:feature:line-models:xnya}.
The parameters of this model are the line slope
($\alpha_1$) and the line intersection on the x-axis ($\alpha_0$).
Nevertheless, this model has also a undetermined form
when $\sum_{i=1}^{n} \left(y_i - \bar{y}\right)^2=0$
(\eg, all the data points with the same $y$ coordinate and different $x$ values).
\begin{equation}\label{eq:feature:line-models:xnya}
x = \alpha_0 + \alpha_1 y
\end{equation}
\begin{equation}\label{eq:feature:line-models:xnya:n}
\hat{\alpha}_1 = \frac{\sum_{i=1}^{n} \left(x_i - \bar{x}\right)\left(y_i - \bar{y}\right)}{\sum_{i=1}^{n} \left(y_i - \bar{y}\right)^2}
\end{equation}
\begin{equation}\label{eq:feature:line-models:xnya:a}
\hat{\alpha}_0 = \bar{x} - \hat{\alpha}_1 \bar{y}
\end{equation}

Instead of only considering a single line model for data fitting,
we select the best line model that fits the laser data.
This selection is based on which model has the lower fitting \gls{rmse},
$\varepsilon_{RMSE}$. \eqref{eq:feature:line-models:ymxb:error} and
\eqref{eq:feature:line-models:xnya:error} represent the fitting \gls{rmse} of
the 2D~line models formulated in \eqref{eq:feature:line-models:ymxb} and
\eqref{eq:feature:line-models:xnya}, respectively.
The validation of the proposed selection approach for the line models is shown
in Figure~\ref{fig:feature:model-error-analysis}, where selecting the model
with the lowest \gls{rmse} achieves the best fitting line estimation.
\begin{equation}\label{eq:feature:line-models:ymxb:error}
\varepsilon_{RMSE,\{\hat{\beta}_0,\hat{\beta}_1\}} = \sqrt{\frac{1}{n} \sum_{i=1}^{n} \left( y_i - \hat{\beta}_0 - \hat{\beta}_1 x_i \right)^2}
\end{equation}
\begin{equation}\label{eq:feature:line-models:xnya:error}
\varepsilon_{RMSE,\{\hat{\alpha}_0,\hat{\alpha}_1\}} = \sqrt{\frac{1}{n} \sum_{i=1}^{n} \left( x_i - \hat{\alpha}_0 - \hat{\alpha}_1 y_i \right)^2}
\end{equation}

\begin{figure}[!t]
\centering
\centerline{\includegraphics[height=0.45\columnwidth]{features/model-error-analysis.eps}}
\caption[Error analysis of simple line regression models evaluating the fitting
\glsentryfull{rmse}]{Error analysis of simple line regression models evaluating
the fitting \glslink{rmse}{\glsentryfull{rmse}}, $\varepsilon_{RMSE}$,
for lines with an inclination angle between -90º and 90º and a simulation step
of 0.25º, a set of 100 data points equally distributed along each line,
a distance Gaussian noise added to the data point coordinates of zero mean
and 0.03~m standard deviation, and 1000 simulation samples per line}
\label{fig:feature:model-error-analysis}
\end{figure}

Other models and fitting algorithms could have been considered to fit a line to
the 2D~laser data, such as formulating a line as the vector to the normal of
the infinite line and the position of the weighted mean of the contributing
points~\cite{pfister:icra:2003}, or using \gls{ransac} to estimate
the line parameters of the generic line model~\cite{fischler-bolles:cacm:1981}.
However, these models and fitting algorithms may require
iterative-based optimization procedures and matrices inversions.
In contrast, our methodology has the advantage of having closed-form equations
for both the line parameters and the fitting error.

\subsection{Feature extraction}\label{sec:features:extraction}

In terms of extracting features from 2D~laser data, first, we need
to transform the scan data from polar to cartesian coordinates.
This transformation preserves the data ordering by its angular coordinate
to speed up the search for the closest consecutive laser points
in the angular direction at a given distance-based neighborhood.

Next, we evaluate each of the laser points $\left(x_i,y_i\right),i=1,\dots,N$
as a possible feature keypoint.
\figref{fig:feature:feature-extraction} presents an example of the clustering
and line fitting results following the proposed methodology on a laser point
$(x_j,y_j)$ at a distance scale $R_k=0.15$~m, where this scale defines
the distance-based neighborhood around the scan point.
These results illustrate the clustering of the arrival
($\left(x_i,y_i\right) \in \{(x_j,y_j)_{\text{arr},k}\}$) and departure
($\left(x_i,y_i\right) \in \{(x_j,y_j)_{\text{dep},k}\}$) neighborhoods,
representing the closest consecutive laser points at the negative and positive
sides of the point $(x_j,y_j)$ in the laser angular direction, respectively.
The intersection point between the fitted lines $line_{\text{arr},k}$ and
$line_{\text{dep},k}$ is considered a candidature feature keypoint
($(x_j,y_j)_{\text{cand},k}$). This intersection point is the result of
the lines intersection depending on which line model
is selected for each neighborhood:
\begin{itemize}
\item $y=\beta_0^{\text{arr},k} + \beta_1^{\text{arr},k} x \enspace$;$\enspace y = \beta_0^{\text{dep},k} + \beta_1^{\text{dep},k} x$
\end{itemize}
\begin{equation}
\left\{
\begin{smallmatrix}
x_{j,\text{cand},k} = \left(\hat{\beta}_{0}^{\text{dep},k}-\hat{\beta}_{0}^{\text{arr},k}\right) / \left(\hat{\beta}_{1}^{\text{arr},k}-\hat{\beta}_{1}^{\text{dep},k}\right)\\
y_{j,\text{cand},k} = \hat{\beta}_{0}^{\text{arr},k} + \hat{\beta}_{1}^{\text{arr},k} \cdot x_{j,\text{cand},k}
\end{smallmatrix}
\right.
\end{equation}
\begin{itemize}
\item $x=\alpha_0^{\text{arr},k} + \alpha_1^{\text{arr},k} y \enspace$;$\enspace y = \beta_0^{\text{dep},k} + \beta_1^{\text{dep},k} x$
\end{itemize}
\begin{equation}
\left\{
\begin{smallmatrix}
y_{j,\text{cand},k} = \left(\hat{\alpha}_{0}^{\text{arr},k}\hat{\beta}_{1}^{\text{dep},k}+\hat{\beta}_{0}^{\text{dep},k} \right) / \left(1-\hat{\alpha}_{1}^{\text{arr},k}\hat{\beta}_{1}^{\text{dep},k} \right)\\
x_{j,\text{cand},k} = \hat{\alpha}_{0}^{\text{arr},k} + \hat{\alpha}_{1}^{\text{arr},k} \cdot y_{j,\text{cand},k}
\end{smallmatrix}
\right.
\end{equation}
\begin{itemize}
\item $y=\beta_0^{\text{arr},k} + \beta_1^{\text{arr},k} x \enspace$;$\enspace x=\alpha_0^{\text{dep},k} + \alpha_1^{\text{dep},k} y$
\end{itemize}
\begin{equation}
\left\{
\begin{smallmatrix}
y_{j,\text{cand},k} = \left(\hat{\alpha}_{0}^{\text{dep},k}\hat{\beta}_{1}^{\text{arr},k}+\hat{\beta}_{0}^{\text{arr},k} \right) / \left(1-\hat{\alpha}_{1}^{\text{dep},k}\hat{\beta}_{1}^{\text{arr},k}\right)\\
x_{j,\text{cand},k} = \hat{\alpha}_{0}^{\text{dep},k} + \hat{\alpha}_{1}^{\text{dep},k} \cdot y_{j,\text{cand},k}
\end{smallmatrix}
\right.
\end{equation}
\begin{itemize}
\item $x=\alpha_0^{\text{arr},k} + \alpha_1^{\text{arr},k} y \enspace$;$\enspace x=\alpha_0^{\text{dep},k} + \alpha_1^{\text{dep},k} y$
\end{itemize}
\begin{equation}
\left\{
\begin{smallmatrix}
y_{j,\text{cand},k} = \left(\hat{\alpha}_{0}^{\text{dep},k}-\hat{\alpha}_{0}^{\text{arr},k}\right) / \left(\hat{\alpha}_{1}^{\text{arr},k}-\hat{\alpha}_{1}^{\text{dep},k}\right)\\
x_{j,\text{cand},k} = \hat{\alpha}_{0}^{\text{arr},k} + \hat{\alpha}_{1}^{\text{arr},k}\cdot y_{j,\text{cand},k}
\end{smallmatrix}
\right.
\end{equation}

\begin{figure}[!t]
\centering
\centerline{\includegraphics[height=0.5\columnwidth]{features/feature-extraction.eps}}
\caption[Example of the proposed corner-like feature extraction methodology
on a laser point $(x_j,y_j)$ at a distance scale $R_k=0.15$~m]{Example of the
proposed corner-like feature extraction methodology on a laser point $(x_j,y_j)$
at a distance scale $R_k=0.15$~m. Scan points $\left(x,y\right)$ simulated
considering a 2D~laser scanner with 0.25º angular resolution,
0.06--10.0~m range, and Gaussian noise with zero mean and 0.03~m
standard deviation (same specifications as the Hokuyo UST-10LX sensor)}
\label{fig:feature:feature-extraction}
\end{figure}

The following indicators determine if a candidate keypoint
is a corner-like feature or not:
minimum of 5 points in each neighborhood,
an absolute angle between arrival and departure fitted lines in the range of
$\left[45\text{º},135\text{º}\right]$,
maximum $\varepsilon_{RMSE}$ in line fitting of 0.025~m,
maximum distance of the lines intersection points on the multiple scales to
the intersection centroid of 0.025~m,
and maximum distance of the intersection point to the scan point
being evaluated of 0.05~m.
The latter indicator avoids having features not representing physical points.
Then, the candidate keypoint is validated if the point passes
all previous indicators in a minimum number of 3 consecutive scales,
to improve robustness to noise on the laser data. The scales considered in the
experiments presented in this chapter are the following ones:
$\left[0.05, 0.075, 0.1, 0.125, 0.15, 0.175, 0.2, 0.225, 0.25\right]$~m.
However, the indicators' thresholds and the scale may be parametrized in
run-time. The feature keypoint definition is
the lines intersection centroid on the multiple scales.

\section{Experimental Results and Discussion}\label{sec:features:results}

TBC

\subsection{Feature extraction results}
\subsection{Features observability}
\subsection{Features association distance}
\subsection{Integration with a line features extractor algorithm}

\section{Summary and Conclusions}\label{sec:features:conclusions}

TBC
